\section{Introduction}\label{sec:intro}
A prepared causal experiment determines which future tests can actually be performed. Their conditional laws determine predictive equivalence. An observer, however, must realize those predictions through one executable program, and the histories on which that program spends its resources are distributed according to the experiment it really runs. These three restrictions---causal realization, actual acquisition, and finite resources---must be kept together when transferring a geometric risk bound from one experiment to another.

There are two obstacles to a long-run theory. First, a finite controller reuses the same report row at every time. A finite-history static reduction does not by itself establish continuity under this reuse. Second, neither bounded mean return times nor a separate result at every discount controls the long-run optimization uniformly. A rare preparation can occupy a substantial fraction of physical time. Its statistical mass and its resource cost must not be removed by conditioning.

We resolve these obstacles for a class defined at the experiment level. A cycle consists of a charged preparation and a causal experiment ending at a visible regeneration boundary. At that boundary the physical preparation is renewed and the observer is forcibly returned to its distinguished state. One stationary $M$-state program is used in every cycle and at every inaccessible model parameter. The reset is part of the protocol, not a consequence attributed to a mixing transition kernel. The within-cycle experiment may be nonlinear, nonmixing, infinite-dimensional, history-dependent, or have changing report supports. The exact theorem assumes domination of each finite report word by a common product reference; it does not assume domination of an infinite path law.

The central result is Theorem~\ref{thm:renewal}. Its hypotheses involve a uniform tail envelope for actual cycle lengths, before a program is optimized. Its conclusions compare the same program class under three resource interfaces: $N$ real acquisition calls, a normalized discount, and average risk. Under a return moment $\sup_{\theta,\gam}\E_\theta^\gam\tau^{1+p}\le K$, $0<p\le1$, they imply
\begin{equation}\label{eq:intro-transfer}
 |V_N^\Theta(M)-V_\av^\Theta(M)|\le 2K N^{-p},\qquad
 |V_\beta^\Theta(M)-V_\av^\Theta(M)|\le 2K(1-\beta)^p.
\end{equation}
No memory is added when the acquisition horizon or discount changes. The bounds concern expected risk uniformly over models and programs; they are not uniform sample-path deviation bounds. Proposition~\ref{prop:time-sharp} shows that the orders in \eqref{eq:intro-transfer} cannot be improved uniformly over the stated moment class.

The acquired geometry is the cycle occupation measure divided by the actual mean cycle length. Theorem~\ref{thm:geometry} bounds online prediction between quantization of that measure and an implementable recursive covering estimate. It uses global covers and stability along actual acquisitions, including all candidate states used by the recursion. Small-ball mass then gives matching curves. Theorems~\ref{thm:hidden-geometry} and \ref{thm:quantum-geometry} obtain their dimension and acquired mass from raw likelihoods and positive instruments, not from a supplied covariance or Fisher matrix. The stability criterion is an occupation-weighted one; a separate heavy-tailed holding class has an exact, possibly singular, memory law without filter contraction.

The resource composition is also sensitive to recurrence. A small total-variation defect on one completed scored cycle need not produce a linear error in time-average risk. Theorem~\ref{thm:defect} gives a tail-dependent modulus, and a $1+p$ moment yields order $\varepsilon^{p/(1+p)}$. Proposition~\ref{prop:joint} realizes this exponent together with a matched memory term and an inaccessible calibration coordinate in one task. Primitive kernel and unnormalized instrument bounds are translated into all-controller moduli in Lemma~\ref{lem:primitive}. Thus the effective Theorem~\ref{thm:compile} does not leave its decisive finite-response error unspecified. Its enumeration is finite, not efficient; table description, report parsing, scratch, and planning cost remain distinct from $M$.

Renewal-reward identities, overshoot estimates, compactness, and finite-intersection arguments are classical tools \cite{Asmussen,Lorden,Kallenberg}. The finite-row tensor argument used here is the one isolated in the discounted common-controller analysis \cite{Discounted}. The new synthesis is the uniform renewal passage at fixed observer size, its sharp defect and time moduli, and its coupling to actual occupation geometry and primitive effective certificates. Section~\ref{sec:comparison} compares this result with strategic-measure existence, average-cost POMDP theory, and finite-memory approximation. In particular, we do not claim that the existence of an average-cost control theory itself is new.
